\documentclass[11pt]{article}
\usepackage[utf8]{inputenc}
\usepackage[T1]{fontenc}
\usepackage{lmodern}
\usepackage{textcomp}
\usepackage{amsmath, amssymb, amsthm}
\usepackage{graphicx}
\usepackage{listings}
\usepackage[hidelinks]{hyperref}

\theoremstyle{plain}
\newtheorem{theorem}{Teorema}[section]
\newtheorem{lemma}[theorem]{Lema}
\newtheorem{proposition}[theorem]{Proposición}
\newtheorem{corollary}[theorem]{Corolario}

\theoremstyle{definition}
\newtheorem{definition}[theorem]{Definición}


\title{c6}
\date{}
\begin{document}
\maketitle

\section{Chapter 6}

\section{Writing Well}

In this chapter we consider some techniques for writing mathematics. We deal with small-scale features: choosing an appropriate terminology and notation, writing clear formulae, mixing words and symbols, writing definitions, introducing a new concept. These are pre-requisites for the more substantial task of structuring and delivering a mathematical argument, to be tackled in later chapters.

Auniversalproblembesetsanyformofspecialisedwriting:howmuchexplanation should we provide? Ideally, we ought to explain the meaning of every word that we use. But this is impossible: we would need to explain the words used in the explanation, and so on. Instead, we should only explain a word or symbol if our explanation will make it clearer than it was before. Accordingly, we shall call a word or symbol primitive1 if it's suitable to use without an explanation of its meaning.

An ordinary English word like `thousand' is obviously primitive, but for more specialised words we must consider the context. When we communicate to the general public, the term multiplication can safely be regarded as primitive. Likewise, there should be no need to explain to a mathematician what an eigenvalue is, while a number theorist will be familiar with conductor. Then there are extremes of specialisation: only a handful of people on this planet will know the meaning of Hsia kernel. Finally, terms such as exceptional set mean different things in different contexts. Understanding what constitutes an appropriate set of primitives is an essential pre-requisite for the communication of complex knowledge. Getting it right is tricky.

The writing developed in this chapter is addressed to a mathematically mature audience. Occasionally, we will consider writing for the general public.

\section{6.1 Choosing Words}

Precision must be the primary concern of anyone who writes mathematics. We list some common mistakes and inaccuracies that result from a poor choice of words.

Bad: the equation x $-$3 $\leq$0 Good: the inequality x $-$3 $\leq$0

1 This terminology is due to Blaise Pascal (French: 1623--1662).

© Springer-Verlag London 2014 F. Vivaldi, Mathematical Writing, Springer Undergraduate Mathematics Series, DOI 10.1007/978-1-4471-6527-9\_6

93

94 6 Writing Well

Bad: the equation x2 $-$1 = (x $-$1)(x + 1) Good: the identity x2 $-$1 = (x $-$1)(x + 1) Bad: the identity x =

$\surd$

x2

$\surd$

Good: the equation x =

x2 Bad: the solution of xk = x Good: a solution of xk = x Bad: the minima of a quadratic function Good: the minimum of a quadratic function Bad: the function sin(x) Good: the sine function Bad: the function f (A) of the set A Good: the image of the set A under f Bad: the interval [1, $\infty$) Good: the ray [1, $\infty$) (the infinite interval [1, $\infty$)) Bad: the set Z minus kZ Good: the set difference of Z and kZ Bad: the area of the unit circle Good: the area of the unit disc Bad: the coordinates of a complex number Good: the real and imaginary parts of a complex number Bad: the absolute value is positive Good: the absolute value is non-negative Bad: the function crosses the vertical axis at a positive point Good: the graph of the function intersects the ordinate axis at a positive point.

Once our writing is accurate, we consider refining the choice of words to differentiate meaning, improve legibility, or avoid repetition.

A set is mostly called a set, although collection or family are useful alternatives (as in the title of Sect.2.3.1). A set becomes a space if it has an added structure, like a metric space (a set with a distance) or a vector space (a set with an addition and a scalar multiplication).

The word element expresses a special kind of subsidiary relationship. If A is a set and a $\in$A, then we say that a is an element of A. The term member is a variant, and so is point, which is the default choice for geometrical sets. So if A $\subset$Rn, then x is a point of A. For sequences, we may replace `element' by term. This is necessary if the elements of a sequence are added (multiplied) together, in which case the operands are the terms of the sum (product), not the elements. However, if V = (vk) is a vector, then vk is a component of V , not a term (even though a vector is just a finite sequence). But if M = (mi, j) is a matrix, then mi, j is an element or an entry of M, not a component (even though a matrix is just a sequence of vectors).

The word variable is used in connection with functions and equations. In the former case it has the same meaning as argument, and it refers to the function's input data. In the latter case it means unknown---a quantity whose value is to be found. Polynomials and rational functions may represent functions or algebraic objects; for the latter, the term indeterminate is preferable to variable.

6.1 Choosing Words 95

The term parameter is used to identify a variable which is assigned a value that remains fixed in the subsequent discussion. For example, the indefinite integral of a function is a function which depends on a parameter, the integration constant:

g(x)dx = f (x) + c.

The two symbols x and c play very different roles here, so we have a one-parameter family of functions, rather than a function of two variables.

A function is not always called a function. This may be true for real functions, but in more general settings the terms map, mapping, and transformation are at least as common. So a complex function C $\to$C may be called a mapping, and a function between euclidean spaces Rn $\to$Rm is a mapping or a transformation. Some functions are called operators. These include basic functions, like addition of numbers or intersection of sets, represented by specialised symbols (+, $\cap$) and syntax (x + y rather than +(x, y)). The arithmetical operators, the set operators, the relational operators, and the logical operators are functions of this type. The attribute binary means that there are two operands, while unary is used for a single operand, like changing the sign of a number. The term operator is also used to characterise functions acting on functions which produce other functions as a result. The differentiation operator is a familiar example. A real-valued function acting on functions is called a functional. So definite integration is a functional.

In logic, a function is called a predicate, a boolean function, or a characteristic function. These terms represent the same thing, but they are not completely interchangeable. `Characteristic function' should be used if there is explicit reference to the associated set, and `predicate' (or `boolean function') otherwise. If the arguments of a function are also boolean, then there is a strong case for using `boolean function'. Some examples:

The negation of a predicate is a predicate. Let $\chi$ be the characteristic function of the prime numbers. Let us consider the boolean function (x, y) $\to$$\neg$x $\land$y. The predicate `n $\to$(3 | n(n + 1))' is the characteristic function of a proper subset of the integers. A relational operator is a boolean function of two variables.

In mathematics it is acceptable to change the meaning of established words. The following re-definitions of universal geometrical terms were taken from mathematics research literature:

By a triangle we mean a metric space of cardinality three. By a segment we mean a maximal subpath of P that contains only light or only heavy edges. By a circle we mean an affinoid isomorphic to max Cp(T, T $-$1).

96 6 Writing Well

Re-defining such common words requires some self-confidence, but in an appropriate context this provocative device may be quite effective. Clearly the new meaning is meant to remain confined to the document in which it's introduced.

\section{6.2 Choosing Symbols}

Choosing mathematical notation is difficult. Mathematicians are notoriously reluctant to accept standardisation of notation, to a degree unknown in other disciplines. Indeed, the ability to adjust quickly to new notation is regarded as one of the skills of the trade. The reality is somewhat different: absorbing new notation requires effort, and most people would gladly avoid it. So if we intend to communicate mathematics without confusing or alienating our audience, the notation must be simple, logical, and consistent.

Two golden rules should inform the use of symbols:

\textbullet{} Do not introduce unnecessary symbols.

\textbullet{} Define every symbol before using it.

Once defined, a symbol should be used consistently. Never use the same symbol for two different things or two symbols for the same thing, even in instances appearing far apart in the document. Don't write `A j, for 1 $\leq$j < n' in one place and `Ak, for 1 $\leq$k < n' in another, unless there is a good reason to do so. Such small inconsistencies produce some `notational pollution'. As the pollution piles up, reading becomes tiresome.

In this section we offer guidelines on how to choose symbols. These are not laws, and may be adapted to one's taste or rejected. What's important is to develop awareness of notation, and to make conscious decisions about it.

Sets. Represent sets by capital letters, Roman or Greek, such as

S A $\Omega$.

The large variety of fonts available in modern typesetting systems increases our choice. When dealing with generic sets, then A, B, C or X, Y, Z are good symbols. For specific sets, choose a symbol that will remind the reader of the nature of the set. So, for an alphabet \{a, b, c, . . .\}, the symbols A, A are obvious choices; likewise,

F is appropriate for a set of functions, etc.

Lower-case symbols like x, y represent the elements of a set. So x $\in$A is a good notation, X $\in$A is bad, and X $\in$a is very bad. If more than one set is involved, consider using matching symbols. Thus

a $\in$A b $\in$B c $\in$C

is more coherent than

x $\in$A y $\in$B z $\in$C.

6.2 Choosing Symbols 97

Some thought may be required for sets of sets. Consider the expressions

f (A $\cap$B) f (x $\cap$y).

The left expression will be interpreted as the image of the intersection of two sets under a function. In this case A$\cap$B represents a subset of the domain of f . However, suppose that the domain of f is a set of sets (e.g., a power set, see Sect.2.3). Then the argument of f is a set, and this expression becomes dangerously ambiguous. The notation in the right expression removes this ambiguity. The combination of standard symbols x, y for variables and a set operator makes it clear that the argument x $\cap$y of the function is an element, rather than a subset, of the domain.

Integers. When choosing a symbol for an integer, begin from the middle region of the Roman alphabet

i, j, k,l, m, n (6.1)

particularly if an integer is used as subscript or superscript.2 However, use p for a prime number, and q if there is a prime different from p. The list of adjacent letters (6.1) cannot be extended; the preceding symbols, f, g, h, are typical function names (see below), while the symbol that follows, o, is rarely used, not only for its resemblance to 0 (zero), but also because it has an established meaning in asymptotic analysis (it appears in expressions of the form o(log x)). A capital letter in the list (6.1) may be used to represent a large integer, or combined with lower-case letters to denote an integer range: n = 1, . . . , N. This combination of symbols is much used in connection with sums and products (Sect.3.2).

Rationals. For rational numbers, use lower-case Roman letters in the ranges a--e, or p--z. The notation

r = m

n

is good, because `r' reminds us of `rational', while numerator and denominator conform to the convention for integers. If there is more than one rational, use adjacent symbols, s, t in this case.

Real numbers. For real numbers, use the same part of the Roman alphabet as for rationals, or the Greek alphabet:

$\alpha$, $\beta$, $\gamma$, . . .

If there are both rational and real numbers and if the distinction between them is important, then use Roman for the rationals and Greek for the reals.

Some Greek symbols have preferential meaning: small quantities are usually represented by $\varepsilon$, $\delta$, while for angles one uses $\varphi$, $\phi$, $\theta$. Famous real constants have dedicated symbols:

2 Physicists use Greek letters for subscripts and superscripts.

98 6 Writing Well

\section{$\pi$ =3.141592653 . . . Archimedes' constant (or Pi) e =2.718281828 . . . Napier's constant3}

(6.2)

\section{$\gamma$ =0.333177924 . . . Euler-Mascheroni constant.4}

These constants are typeset with the `upright' typeface, to highlight their distinguished status. Upright fonts are known as roman (even for Greek letters), as opposed to italic, which are slanted fonts: $\pi$, e, $\gamma$ .

Complex numbers. Complex numbers tend to occupy the end of the Roman alphabet, and your first choice should be z or w. On the complex plane, we write z = x+iy, typeset in roman font. (However, number theorists use $\surd$$-$1, not i.) The standard where x and y are the real and imaginary parts of z, and i is the imaginary unit, again

notation for a complex number in polar coordinates is z = $\rho$ei$\theta$, which also combines roman and italic fonts.

Unknowns. The quintessential symbol for an equation's unknown is x, invariably followed by y and z if there are more unknowns. For a large number of unknowns it is necessary to use sequence notation x1, . . . , xn. This notation is also appropriate for the indeterminates of a polynomial and for the arguments of a function.

Composite objects. Objects which have constituent elements (groups, graphs, matrices) are best represented with capital letters, Roman or Greek. So use G or $\Gamma$ for a group or a graph, and M for a matrix. If you have two groups, use adjacent symbols, like G and H. As with sets, for these objects' components consider using matching symbols, e.g., g $\in$G. A notable exception are graphs, where v and e are invariably used for vertices and edges, respectively.

Functions. The default choice for a function's name is, of course, f , and if there is more than one function, use the adjacent symbols g, h. Lower-case symbols work well with any number of variables: f (x), f (x, y, z), f (x1, . . . , xn). But if the co-domain of a function is a cartesian product, then the function is a vector, and capital letters are preferable. So a real function of two variables may be specified as

F: R2 $\to$R2 (x, y) $\to$( f1(x, y), f2(x, y)).

Greeksymbols,eithercapitalorlower-case,arealsocommonlyusedforfunctions' names. The contrast between Roman and Greek symbols may be exploited to separate out the symbols' roles, as in $\mu$(x) or f ($\lambda$).

Some famous functions are named after, and represented by, a symbol (often a Greek one), thereby creating a strong bond between object and notation. The best known are Euler's gamma function

$\infty$

xs$-$1e$-$xdx

(s) =

0

3 John Napier (Scott: 1550--1617). 4 LorenzoMascheroni (Italian: 1750--1800).

6.2 Choosing Symbols 99

(the extension of the factorial function to complex arguments), and Riemann's zetafunction5 $\zeta$

$\infty$

1 ns . (6.3)

$\zeta$(s) =

n=1

There is a peculiar notation for this function: its complex argument s is commonly written as s = $\sigma$ + i$\tau$, with $\sigma$ and $\tau$ real numbers. Other functions with dedicated notation are Euler's $\phi$-function (Eq. 3.8), Dedekind's $\eta$-function, Kroneker's $\delta$-function, Weierstrass' P-function, Lambert's W -function, etc.

Sequences and vectors. Sequences pose specific notational problems due to the presence of indices. Consider the various possibilities listed in (3.2): which one should we choose? Be guided by the principle of economy: a symbol should be introduced only if it's strictly necessary. So the notation (ak) is quite adequate for a generic sequence, or if the specific properties of the sequence (the initial value of the index, its finiteness) are not relevant. When more information is needed, the notation (ak)k$\geq$1 is more economical than (ak)$\infty$

k=1, but the latter may be a better choice if it is to be contrasted with (ak)n

k=1. In turn, the latter is not as friendly as (a1, . . . , an), although it is more concise.

If a sequence is referred to often, even the stripped down notation (ak) could become heavy, and it may be advisable to allocate a symbol for the sequence:

a = (a1, a2, . . .) v = (v1, . . . , vn).

As usual, we have employed matching symbols, using, respectively, a minimalist lower-case Roman character and a lower-case boldface character which is common for vectors. When using ellipses, two or three terms of the sequence usually suffice, but there are circumstances where more terms or a different arrangement of terms is needed.

For example, in the expression

(1 + x, 1 + x2, . . . , 1 + x2k, . . .)

the insertion of the general term removes any ambiguity, while the ellipsis on the right suggests that the sequence is infinite---cf. (3.1). The notation

(a1, . . . , ak$-$1, ak+1, . . . , an)

denotes a sub-sequence of a finite sequence obtained by deleting the k-th term, for an unspecified value of k = 1, n.

Things get complicated with sequences of sequences. This situation is not at all unusual; for instance, we may have a sequence of vectors whose components must be referred to explicitly. We write

5 Bernhard Riemann (German: 1826--1866).

100 6 Writing Well

V = (V1, V2, . . .) or v = (v1, v2, . . .).

Let Vk (or vk) be the general term of our sequence. How are we to represent its components? As usual, we choose the matching symbol v, with a subscript indicating the component. However, the integer k has to appear somewhere, and its range must be specified. It is advisable to keep k out of the way as much as possible:

Vk = (v(k)

1 , . . . , v(k) n ) k = 1, 2, \dots{}

In this expression we have used ellipses to specify the ranges of the indices; we could have used inequalities as well:

Vk = (v(k)

j ) 1 $\leq$j $\leq$n, k $\geq$1.

The parentheses are obviously needed for the superscript, for otherwise vk

i would be interpreted as vi raised to the k-th power. However, it may just happen that we need to raise the vector components to some power. Clearly we can't use adjacent superscripts v(2) 4

, so parentheses are needed, but the straightforward notation (v(2)

3 )4

3

is awkward. For a more elegant solution, we represent k as an additional subscript, adopting, in effect, matrix notation:

Vk = (v1,k, . . . , vn,k) k $\geq$1.

As a side note, one should keep in mind that with vectors the multiplication symbols `\textperiodcentered{}' and `$\times$' are reserved for the scalar and vector products, respectively---cf. (2.10). Hence for scalar multiplication we must use juxtaposition:

a(bV \textperiodcentered{} cW) x v $\times$ y u.

Derived symbols. Closely related objects require closely related notation. Proximity in the alphabet, e.g., x, y, z may be used for this purpose. For a stronger bond, the meaning of a symbol may be modified using subscripts, superscripts and other decorations:

A$*$ $\eta$ n+ h \textasciitilde{}e $\Omega$$-$ Zr.

The derived symbol N0 = N $\cup$\{0\} is often found in the literature. Many symbols derived from R are in use:

R+ = \{x $\in$R : x > 0\} R$\geq$0 = \{x $\in$R : x $\geq$0\}

R$*$= R$\setminus$\{0\} R = R $\cup$\{$\infty$\}.

6.2 Choosing Symbols 101

These sentences illustrate the use of derived symbols:

Let f :X $\to$X be a function, and let x$*$= f (x$*$) be a fixed point of f . We consider the endpoints x$-$and x+ of an interval containing x. Let f be a real function, and let

f (x) if f (x) $\geq$0 0 if f (x) < 0.

f + : R $\to$R x $\to$

It must be noted that there is no general agreement on the meaning of decorations. Thus for a set the over-bar denotes the so-called closure---adjoining to a set all its boundary points, see Sect.5.4.1 (the transformation from R to R is essentially a closure operation). However, for complex numbers the over-bar denotes complex conjugation. If f is a function, then f $'$ is the derivative of f , but for sets a prime indicates taking the complement.

Example. We illustrate notational problems raised by the coexistence of variables and parameters. For a fixed value of z, the bivariate polynomial

f (x, z) = $-$z2 + xz + 1

becomes a polynomial in x, and we wish to adopt a notation that reflects the different roles played by the symbols x and z. We replace z with a, to keep it far apart from x in the alphabet, and we rewrite the expression above as

fa(x) = ax + 1 $-$a2 a $\in$R. (6.4)

We now have a one-parameter family of linear polynomials in x. For fixed a, the equation y = fa(x) is the cartesian equation of a line, so we also have a one-parameter family of lines in the plane. Plotting some of these lines reveals a hidden structure: they form the envelope of a parabola (Fig. 6.1). Likewise, if we fix x = a we obtain a one-parameter family of quadratic polynomials ga(z) = $-$z2 + az + 1.

\section{6.3 Improving Formulae}

When the physicist Stephen Hawking was writing his book A Brief History of Time, an editor warned him that every new equation in the book would halve the readership. So he included a single equation. When writing for the general public, one should expect diffidence---even hostility---towards formulae. Even though mathematicians are trained to deal with formulae, it is still safe to assume that no one likes to struggle with too many symbols. In this section we explore techniques to improve the clarity of formulae, through presentation, notation, and layout.

102 6 Writing Well

Fig. 6.1 The one-parameter family of lines y = fa(x), where fa is given by (6.4). These lines are tangent to the parabola y = x2/4 + 1

4

3

2

1

0

-4 -2 0 2 4

x

-1

Formulae may be embedded in the text or displayed, and in a document one usually finds both arrangements. In either case a formula must obey standard punctuation rules. The following passage features embedded formulae with appropriate punctuation.

For each x $\in$X we have the decomposition x = $\xi$ +$\lambda$, with $\xi$ $\in$ and $\lambda$ $\in$$\Lambda$; accordingly, we define the function P : X $\to$$\Xi$, x $\to$$\xi$, which extracts the first component of x.

The punctuation generates rests as it would in an ordinary English sentence. The two items comprising the function definition are separated by a comma, which would not be necessary in a displayed formula, see (2.16). The formula defining P is echoed by a short sentence, for added clarity.

This is a fully punctuated displayed formula:

a2

k $-$1, 1 $\leq$k < 10;

a1 = 1; ak+1 =

a2

k, k $\geq$10.

If full punctuation seems heavy (it may do so here), then we may replace some punctuation by increased spacing or by words:

a2

k $-$1 if 1 $\leq$k < 10; a2

a1 = 1; ak+1 =

if k $\geq$10.

k

A displayed formula is normally embedded within a sentence, and the punctuation at the end of a formula must be appropriate to the structure of the sentence. In particular, if a sentence terminates at a formula---as in the example above---then the full stop at the end of the formula must always be present.

6.3 Improving Formulae 103

Example. The following untidy formula

f (x) = 14x $-$2x3 $-$2x2 + 14

$-$2x $-$4

couldrepresentatypicalunprocessedoutputofacomputeralgebrasystem.Itcontains redundant information (a common factor between numerator and denominator), the monomials at numerator are not ordered, and there are too many negative signs. The properties of the rational function f (x) are not evident from it. There are several ways to improve the formula's layout:

f (x) = x3 + x2 $-$7x $-$7

x + 2

f (x) = (x + 1)(x2 $-$7)

x + 2

f (x) = x2 $-$x $-$5 + x + 2. 3

If the degree or the coefficients of f are important, then the first version is appropriate; the second version makes it easy to solve the equation f (x) = 0; the third version is a preparation for integrating f .

Example. The following subset of the rationals

x y $\in$Q : y = x2 + 1, x $\in$Z, x < 0

A =

is defined with a disproportionate volume of notation. To save symbols we switch from the Zermelo to the standard set definition, and consider only elements of the required form. We also remove the inequality by introducing a negative sign, so we can replace the integers by the natural numbers:

$-$n n2 + 1 : n $\in$N

.

A =

Now the formula is transparent.

Example. To simplify a complex formula, we apply the `divide and conquer' method. In the cluttered expression

R = x(ad $-$bc) $-$y(ad $-$cb)2 + z(ad $-$cb)3

thesub-expressionad$-$bc appearsasaunit.Weexploitthisfacttoimprovethelayout:

R = x$\Delta$ $-$y$\Delta$2 + z$\Delta$3 $\Delta$ = ad $-$bc.

The formula is tidier, and the structure of R is clearer.

104 6 Writing Well

Example. Our next challenge is to improve an intricate double sum:

yi$-$1

$\infty$

(y + 1)xi$-$1. (6.5)

z(y1, y2, . . .) =

y=0

i=1

The meaning of z is not at all evident. This quantity depends on x, but its dependence does not appear explicitly; a poor choice of symbols obscures matters further. We note that the parameters yi are integers, since they are the upper limit of the inner summation; accordingly, we replace the symbol y with n---see (6.1). Then we adopt the `divide and conquer' method, splitting up the sum as follows:

ni$-$1

$\infty$

di xi$-$1 di =

z(x, n1, n2, . . .) =

(k + 1).

i=1

k=0

We now see that z is a power series in x; its coefficients are finite sums, determined by the elements of an integer sequence. These are sums of arithmetic progressions, which can be evaluated explicitly:

k = n(n + 1)

. 2

(k + 1) =

0$\leq$k$\leq$n$-$1

1$\leq$k$\leq$n

(In this passage we have dropped the subscript i, since the association n $\leftrightarrow$ni is clear, and we have switched to the standard sigma-notation to change variable--- cf. Eq. (3.7)). Our original double sum (6.5) can now be written as

$\infty$

z(x, n) = 1

ni(ni + 1)xi$-$1 n = (n1, n2, . . .)

2

i=1

where we have used again the `divide and conquer' method. The dependence of z on the variable x and the sequence n is now clear.

\section{6.4 Writing Definitions}

Adefinitionrequiresapause,togivethereadertimetoabsorbit.Thismaybeachieved by giving the definition twice, first with words then with symbols (or vice-versa), by using two different formulations, or by supporting the definition with an example. Let us analyse some definitions.

Example.

Let P be a predicate over the rational numbers, that is, a function of the type

P : Q $\to$\{T, F\}.

6.4 Writing Definitions 105

The definition of the symbol P uses some jargon (predicate over a set) so the second part of the sentence recalls its meaning. A previous exposure to this concept is tacitly assumed; for a first encounter we would need a more considerate style, such as in the opening paragraph of Sect.4.3.

Example.

For given $\varepsilon$ > 0, we define the set

$\surd$

S$\varepsilon$ = \{(x, y) $\in$R2 : |x $-$y| $\leq$$\varepsilon$/

2\},

that is, S$\varepsilon$ is a strip of width $\varepsilon$ symmetrical with respect to the main diagonal in the cartesian plane.

This time the symbolic (Zermelo) definition appears first, while a verbal explanation clarifies the geometric meaning of S$\varepsilon$, which is not obvious from the formula. The quantity $\varepsilon$ appears as a subscript, indicating that it is a parameter. Zermelo definitions should be used sparingly, and should not be used at all if we write for a non-mathematical audience. (Lars Ahlfors in his beautifully written text Complex analysis deliberately avoids them [2].)

Example.

For every real number $\lambda$, let $\Pi$($\lambda$) be the plane in three-dimensional euclidean space orthogonal to the vector v($\lambda$) = (1, $\lambda$, $\lambda$2). Thus $\Pi$($\lambda$) consists of the points z = (x1, x2, x3) $\in$R3 for which the scalar product z \textperiodcentered{} v = x1 + x2$\lambda$ + x3$\lambda$2 is equal to zero.

In this definition the quantity $\lambda$ appears as an argument of both $\Pi$ and v, and thus the latter represent functions. The second sentence restates the definition in a form suitable to computation, recalling the connection between orthogonality and scalar product, and introducing further notation.

Example.

Let N be the set of sequences of natural numbers, such that every natural number is listed infinitely often. For example, the sequence

(1, 1, 2, 1, 2, 3, 1, 2, 3, 4, . . .)

belongs to N .

A non-experienced reader will have no idea that such a construction is at all possible, so we give an example.

The definition of a symbol should appear as near as possible to where the symbol is first used; defining a symbol immediately after its first appearance is also acceptable, provided that the definition is given within the same sentence.

Example. We give the same definition three times, articulating the changes in emphasis that accompany each version.

106 6 Writing Well

Consider the power series

$\infty$

anxn,

h(x) =

n=1

where the coefficient an is the square of the n-th triangular number.

The definition of an immediately follows its first appearance. The displayed formula represents a general power series, and its specific nature is revealed only by reading the entire sentence. For this reason, this format may not be ideal if the formula is to be referred to from elsewhere in the text. This definition puts some burden on readers who are unfamiliar with the term triangular number.

In our second version we change both text and formula.

Let tn be the n-th triangular number. We consider the power series

$\infty$

n xn.

t2

h(x) =

n=1

Now tn is defined before being used, with the symbol t chosen so as to remind us of `triangular'. The formula is clearer. Just glancing at it makes us want to find out what tn is, and to do this one would begin to scan the text preceding, rather than following, the formula.

Our third version combines verbal and symbolic definitions.

We consider the power series h(x) whose coefficients are the square of the triangular numbers, namely,

$\infty$

n xn tn = n(n + 1)

t2

. 2

h(x) =

n=1

The formula is more complex because everything is defined within it. But it is also self-contained, so it'll be easy to refer to.

\section{6.5 Introducing a Concept}

Opening the exposition with a formal definition is rarely a good way to introduce a new concept, particularly if we write for non-experts. The following examples show how to set the scene for a gentle introduction to a new idea. In each case we ask a question---a simple rhetorical device to engage and prepare an audience.

6.5 Introducing a Concept 107

Example. Introducing recursive sequences.

Let n be a natural number. How is the power 2n defined? We could use repeated multiplication

2n := 2 $\times$ 2 $\times$ \textperiodcentered{} \textperiodcentered{} \textperiodcentered{} $\times$ 2

,

n

but we could also write

21 := 2 and 2n := 2 $\times$ 2n$-$1 n > 1.

The second formula is an example of a recursive definition of a sequence (an). When n = 1, the first term a1 = 2 is defined explicitly; then, assuming that an$-$1 has been defined, we define an in terms of it. [The recursive definition of a general sequence follows (see Sect.9.4)].

The opening question leads to the familiar definition of exponentiation in terms of multiplication. The less familiar recursive definition comes after, illustrated by a simple example and supported by an appropriate notation. After this preamble, the reader should be ready to confront an abstract definition.

Example. Introducing the exponential function.

The process of differentiation turns a real function into another real function. For example, differentiation turns the sine into the cosine. Are there functions that are not changed at all by differentiation?

[The definition of the exponential function follows.]

We recall a structural property of differentiation and present a familiar example. The question that follows suggests how the argument will develop.

Example. Introducing the rational numbers.

A rational number is represented by a pair of integers, the numerator and the denominator. As these integers need not be co-prime, we may choose them in infinitely many ways. How are we to construct a single rational number from an infinite set of pairs of integers? How do we define the set Q from the cartesian product Z $\times$ (Z$\setminus$\{0\})? [The definition of the set of rational numbers follows.]

The first two sentences recall elementary facts. A question then invites the reader to think about this problem more carefully. A second question, which echoes the first, uses proper terminology and notation, in preparation for a formal construction.

Example. We design an exercise structured as a list of questions. The topic is the number of relations on a finite set.

108 6 Writing Well

What is a relation on a set? Can a relation be defined on the empty set? How many relations can one define on a two-element set? Let n be the number of relations on a set. What values can n assume?

The first question checks background knowledge; the other questions gently explore the problem, guiding the reader from the specific to the general. Textbook exercises are sometimes structured in this way, to encourage independence in learning.

Formulating questions is not just a method to grab attention or structure exercises. Asking the right questions---those that chart the boundaries of our knowledge---is the essence of research.

\section{6.6 Writing a Short Description}

Writing a synopsis of a mathematical topic is a common task. It could be the abstract of a presentation, the summary of a chapter of a book, an informal explanation of a theorem. It could be the closing paragraph of a large document, which distils the essence of an entire subject.

Writing a short essay is difficult---the shorter the essay, the greater the difficulty. We shall adopt the format of a micro-essay: 100--150 words (one or two paragraphs) and no mathematical symbols. Within such a confined space one is forced to make difficult decisions on what to say and what to leave out; the lack of symbols gives further prominence to the concepts. Our command of the syntax will be put to the test.

Our first micro-essay is a summary of Sect.2.1 on sets. We have access to all relevant material, and the main difficulty is to decide what are the highlights of that section. We select two ideas: how to define a set, and how to construct new sets from old ones.

A set is a collection---finite or infinite---of distinct mathematical objects, whose ordering is unimportant. A small set may be defined by listing all its elements explicitly; a large set is instead defined by specifying the characteristic properties of the elements.

Combining numbers with arithmetical operators gives new numbers. Likewise, combiningsetsviasetoperators(union,intersection,difference)givesnewsets. But this is a bit like recycling what we already have. A class of brand-new sets are the so-called cartesian products, which are constructed by pairing together elements of existing sets. A well-known example is the cartesian plane, which is made of pairs of elements of the real line, each representing a coordinate.

The description of set operators exploits an analogy with arithmetical operators, and for this reason the first two sentences of the second paragraph have the same structure. We avoid a precise definition of the term `cartesian product' (not enough space!), opting instead for an informal description and an illustrative example.

6.6 Writing a Short Description 109

Next we write a micro-essay on prime numbers, a synthesis of our knowledge of this topic.

A prime is a positive integer divisible only by itself and unity. (However, 1 is not considered prime). The importance of primes in arithmetic stems from the fact that every integer admits a unique decomposition into primes. The infinitude of primes---known from antiquity---and their unpredictability make them object of great mathematical interest. Many properties of an integer follow at once from its prime factorization. For instance, looking at the exponents alone, one can determine the number of divisors, or decide if an integer is a power (i.e.,6 a square or a cube).

Primality testing and prime decomposition are computationally challenging problems with applications in digital data processing.

This essay begins with a definition. The technical point concerning the primality of 1 has been confined within parentheses, to avoid cluttering the opening sentence. The following two sentences deliver core information in a casual---yet precise---way. We state two important theorems (the Fundamental Theorem of Arithmetic, and Euclid's theorem on the infinitude of the primes) without mentioning the word `theorem'. We also give a hint of why mathematicians are so fascinated by primes. The second paragraph elaborates on the importance of unique prime factorisation, by mentioning two applications without unnecessary details. The short closing paragraph, like the last sentence of the first paragraph, is an advertisement of the subject, meant to encourage the reader to learn more.

Now a real challenge: write a micro-essay on Theorem 4.2, which we reproduce here for convenience.

Theorem. Let X be a set, let A, B $\subseteq$X, and let PA, PB be the corresponding characteristic functions. The following holds (the prime denotes taking complement):

(i) $\neg$ PA = PA$'$

(ii) PA $\land$ PB = PA $\cap$B

(iii) PA $\lor$ PB = PA $\cup$B

(iv) PA $\Rightarrow$PB = P(A$\setminus$B)$'$

(v) PA $\Leftarrow$$\Rightarrow$PB = P(A$\cap$B) $\cup$(A $\cup$B)$'$.

We must extract a theme from a daunting list of inscrutable symbols. We inspect the material of Sect.4.3 leading to the statement of this theorem: it deals with characteristic functions, and the main idea is to link characteristic functions to sets. We can see such a link in every formula: on the left there are logical operators, on the right set operators. Given the specialised nature of this theorem, some jargon is unavoidable, so we write for a mathematically mature audience.

Every characteristic function corresponds to a set, and vice-versa. If we let a logical operator act on characteristic functions we obtain a new characteristic

6 Abbreviation for the Latin id est, meaning `that is'.

110 6 Writing Well

function, and with it a new set. How is this set related to the original sets? This theorem tells us that for each logical operator there is a set operator which mirrors its action on the corresponding sets. In other words, we have a biunique correspondence between these two classes of operators. For instance, the logical operator of negation is represented by the set operation of taking the complement. More precisely, the negation of the characteristic function of a set is the characteristic function of the complement of this set. Analogous results are established with respect to the other logical operators.

Here words are better than symbols to describe the structure, but they can't compete with symbols for the details. So, to avoid tedious repetitions, we have chosen to explain just one formula carefully (the easiest one!), mentioning the other formulae under the generic heading `analogous results'.

Exercise 6.1 Answer concisely. [ $\varepsilon$, 30]

1. What is the difference between an ordered pair and a set?

2. What is the difference between an equation and an identity?

3. What is the difference between a function and its graph?

Exercise 6.2 Consider the following question:

Why is it that when the price of petrol goes up by 10 \% and then comes down 10 \%, it doesn't finish up where it started?

1. Write an explanation for the general public. Do not use mathematical symbols, as most people find them difficult to understand. [ $\varepsilon$ ]

2. Write an explanation for mathematicians, combining words and symbols for maximum clarity. You should deal with the more general problem of two opposite percentage variations of an arbitrary non-negative quantity, not specialised to petrol, or to 10\% variation, or to the fact that the decrease followed the increase and not the other way around.

Exercise 6.3 Consider the following question:

I drive ten miles at 30 miles an hour, and then another ten miles at 50 miles an hour. It seems to me my average speed over the journey should be 40 miles an hour, but it doesn't work out that way. Why not?

Write an explanation for the general public, clarifying why such a confusion may arise. You may perform some basic arithmetic, but do not use symbols as most people find them difficult to understand. [ $\varepsilon$ ]

Exercise 6.4 Consider the following question:

I tossed a coin four times, and got heads four times. It seems to me that if I toss it again I am much more likely to get tails than heads, but it doesn't work out that way. Why not?

6.6 Writing a Short Description 111

Write an explanation for the general public, clarifying why such a confusion may arise. You may use symbols such as H and T for heads-tails outcomes, but avoid using other symbols. [ $\varepsilon$ ]

Exercise 6.5 Consider the following question:

In a game of chance there are three boxes: two are empty, one contains money. I am asked to choose a box by placing my hand over it; if the money is in that box, I win it. Once I have made my choice, the presenter---who knows where the money is---opens an empty box and then gives me the option to reconsider. I can change box if I wish. It seems to me that changing box would make no difference to my chances of winning, but it does not work out that way. Why not?

Write for the general public, explaining what is the best winning strategy. [ $\varepsilon$ ]

Exercise 6.6 Write a micro-essay on each topic. [ $\varepsilon$, 150]

1. Quadratic equations and complex numbers.

2. Definite versus indefinite integration.

3. What are matrices useful for?

4. Images and inverse images of sets.

5. Rational numbers versus fractions.

6. How to get the integers from pairs of natural numbers.

Exercise 6.7 Write a two-page essay on each topic.

1. The many ways of defining the logarithm.

2. All fractions in a very small interval have very large denominators (with at most

one exception, that is).

3. In probability, a random variable is not random and is not a variable: it's a

function!

4. Sequences racing to infinity: who gets there first?

\end{document}